\section{与人协作}
\label{sec:person}

基准测试只把一个完整的任务交给智能体一次；而一个真实的人会纠正它、过后再回来，并期望它记得。四个产品测试台测量了
这一面（表~\ref{tab:benches}）。它们的结果与编程实验指向同一个方向：能传下去的，是写在下一个智能体会读到的地方的
东西，而不是存在记忆工具里的东西。

\begin{table}[t]
\centering
\footnotesize
\caption{产品与人协作的测试台。由脚本扮演的人和小组代替真人；“之前”和“之后”指最后一列所列改动之前和之后的产品。}
\label{tab:benches}
\begin{tabularx}{\textwidth}{@{}>{\raggedright\arraybackslash}p{20mm} >{\raggedright\arraybackslash}p{38mm} L >{\raggedright\arraybackslash}p{24mm}@{}}
\toprule
测试台 & 设置 & 结果 & 上线的改动 \\
\midrule
验收习惯（11/17--18） & 一个脚本扮演的人，有四条任何任务都没写明的代码习惯；每个序列六道 Python 题；过程中给出纠正 &
  之前：一条被纠正过的习惯，经由主助手时 51 次中 0 次再被违反（它把这个人的原话贴进每份任务简报），直接交给团队时
  51 次中 51 次被违反。在留出的人上：团队违反被纠正习惯 14/96，主助手 0/48；一次性要求被误当成规则保留 0/108 和
  3/81 & 常设规则文件和一个保存规则的工具 \\
记忆教训（11/18） & 29 个序列，后面的任务需要用到前面任务的教训 & 全新运行 0/6，告知后 4/4，团队 6/8，带记忆的团队 5/8，
  主助手 8/8；团队恰好在读了前一任务目录时通过（11/11 对 0/5）；压缩之后 12/12 & 不重写记忆；教训通过文件和任务简报
  传递 \\
跟进（10/25） & 脚本化的小组把 7 种回测变体之一交给主助手；每组 21 次运行 & 跟进率前后都是 18/21；评分细则得分从 0.49
  升到 0.88 & 主助手的交回文字 \\
决策模型（10/25--26） & 盲审阅读者判断一个小组是否揭示了脚本设定的决策模型 & 有实验设计技能时 24/24 对 16/24
  （$p = 0.004$），token 为 0.88 倍 & 实验设计技能 \\
\bottomrule
\end{tabularx}
\end{table}

有两个结果改变了设计。主助手之所以守住每条习惯，是因为它把这个人的原话改写进了每份任务简报；团队只有在这些话进入
它们会读的文件时才守得住。并且在 104 个序列任务中，没有一个智能体主动调用过记忆工具；所以常设规则现在放在下一个
智能体会读的文件里，教训则留给工作本身的文件和主助手的任务简报去传递。

有一条工作线没有找到值得上线的东西。“进化实验室”跑了八轮（模型成本 59.55 美元），追查一次失败，结果发现那次失败
主要来自我们自己写的任务摘要：摘要改正之后，上线团队 4 次中 4 次解出该题，之前是 21 次中 3 次。
